\chapter{Apprendimento non supervisionato: K-means e SOM}
\label{cap:nonsup}

\begin{inbreve}
Il capitolo introduce l'apprendimento non supervisionato dal punto di vista del clustering come quantizzazione
vettoriale. Si ricava il K-means (versione on-line e batch) come discesa del gradiente sull'errore di quantizzazione e
si arriva alle Self-Organizing Map di Kohonen, che al clustering aggiungono la preservazione della topologia e quindi
la visualizzazione dei dati.
\end{inbreve}

\section{Apprendimento non supervisionato}

Nell'apprendimento non supervisionato il training set contiene solo dati \textbf{non etichettati} $\langle\x\rangle$
(\rif{sec:task}). I task principali sono il \textbf{clustering}, la \textbf{riduzione di dimensionalità} (con
visualizzazione e pre-elaborazione: PCA, Multi-Dimensional Scaling, Independent Component Analysis, \dots) e la
modellazione della densità (non trattata). Un motivo pratico della sua importanza è che i \textbf{dati non etichettati
costano molto meno} di quelli etichettati.

\begin{definizione}[Clustering]
Il \termine{clustering} partiziona i dati in \textbf{cluster}, sottoinsiemi di dati simili: gli esempi di un cluster
sono più simili tra loro che agli esempi di un altro cluster. Ogni cluster è rappresentato da un \termine{centroide}
(prototipo, vettore di riferimento, \ingl{codevector}); l'insieme dei centroidi è il \termine{codebook}.
\end{definizione}

\begin{attenzione}[La metrica domina]
Come nel K-NN, è la \textbf{distanza assunta} a stabilire quali esempi sono simili, e i raggruppamenti dipendono da
quella metrica. La distanza euclidea è popolare ma è \textbf{solo una scelta}, con implicazioni profonde (ad esempio si
possono pesare di meno gli outlier usando $1$ meno una gaussiana centrata sul prototipo).
\end{attenzione}

\section{Quantizzazione vettoriale}

\begin{definizione}[Vector Quantization]
La \termine{quantizzazione vettoriale} (VQ) codifica una varietà di dati $V\subseteq\R^n$ usando un insieme finito
di vettori di riferimento (il codebook) $\w = (\w_1,\dots,\w_K)$, $\w_i\in\R^n$. Un vettore $\x\in V$ è descritto dal
vettore di riferimento \textbf{vincente} (\ingl{best matching}) $\w_{i^*(\x)}$, quello con distorsione
$d(\x,\w_{i^*(\x)})$ minima.
\end{definizione}

La VQ divide la varietà nei \textbf{poliedri di Voronoi} (\rif{sec:knn-algoritmo}, \figrif{fig:voronoi})
\[
V_i = \{\x\in V : \norm{\x-\w_i}\le\norm{\x-\w_j}\ \forall j\},
\]
e ogni $\x$ è descritto dal vettore di riferimento della sua cella. Il problema è l'inverso di quello del K-NN:
\textbf{dato il dataset, trovare i centri} che minimizzano la loss. In generale è difficile: trovare il codebook ottimo è
\textbf{NP-completo}.

\begin{esempio}[Quantizzazione in una dimensione]
Nell'elaborazione dei segnali la quantizzazione approssima un intervallo continuo di valori con un piccolo insieme di
simboli discreti, commettendo un errore di quantizzazione (distorsione): ogni valore è sostituito dal centroide del suo
intervallo.
\begin{center}
\begin{tikzpicture}[x=1.2cm]
\draw[->] (0,0) -- (9.3,0);
\foreach \t in {1.5,3.2,4.6,6.4,8.0} \draw (\t,-0.15) -- (\t,0.15);
\foreach \c in {0.75,2.35,3.9,5.5,7.2,8.65} \fill[blu] (\c,0) circle (2.2pt);
\node[font=\scriptsize, grigio] at (2.35,-0.4) {cella};
\node[font=\scriptsize, blu] at (2.35,0.35) {centroide};
\end{tikzpicture}
\end{center}
\end{esempio}

Il clustering ha un obiettivo più generale (trovare raggruppamenti \emph{interessanti}, spesso definiti implicitamente
dalla procedura stessa e non dall'errore di quantizzazione minimo), ma la VQ fornisce un quadro utile e i suoi algoritmi
sono molto usati per il clustering partizionale.

\subsection{L'errore di quantizzazione}

L'obiettivo è partizionare in modo ottimo una distribuzione sconosciuta in regioni approssimate da prototipi, cioè
trovare un quantizzatore $\x\mapsto c(\x) = \w_{i^*(\x)}$ da uno spazio continuo a uno discreto. Con la distorsione
quadratica $d(\x,c(\x)) = \norm{\x-c(\x)}^2$ (\rif{sec:loss}), la loss è la distorsione attesa (errore di
ricostruzione o di quantizzazione):
\[
E = \int\norm{\x-\w_{i^*(\x)}}^2\,p(\x)\,d\x ,
\]
dove $p(\x)$ è la distribuzione degli input. In versione discreta, su $l$ dati e $K$ prototipi,
\begin{equation}\label{eq:errore-quantizzazione}
E = \sum_{i=1}^{l}\sum_{j=1}^{K}\norm{\x_i-\w_j}^2\,\delta_{\mathit{winner}}(i,j),\qquad
\delta_{\mathit{winner}}(i,j) = \begin{cases}1 & \text{se } j \text{ è il vincitore per } \x_i\\ 0 & \text{altrimenti.}\end{cases}
\end{equation}
Minimizzare $E$ rispetto ai $\w$ è il problema della VQ. L'integrando \textbf{non è differenziabile con continuità},
perché il vincitore cambia in modo discreto al variare di $\x$; le derivate si calcolano però \textbf{localmente},
fissando $\x$ dentro una cella di Voronoi.

\section{K-means}

\subsection{Versione on-line}

\begin{proprieta}{Regola del K-means on-line}
La discesa del gradiente stocastica su \eqref{eq:errore-quantizzazione} dà, per ogni esempio $\x_i$,
\[
\Delta\w_{i^*} = \eta\,\delta_{\mathit{winner}}(i,i^*)\,(\x_i-\w_{i^*}) :
\]
si sposta \textbf{solo il prototipo vincitore}, verso $\x_i$.
\end{proprieta}
\begin{proof}
Per un esempio $\x_i$ fissato, il termine di $E$ che dipende da $\w_j$ è $\norm{\x_i-\w_j}^2\delta_{\mathit{winner}}(i,j)$.
Poiché $\norm{\x_i-\w_j}^2 = \sum_t(x_{i,t}-w_{j,t})^2$, la derivata rispetto a $\w_j$ è
$-2(\x_i-\w_j)\,\delta_{\mathit{winner}}(i,j)$: è nulla per tutti i prototipi tranne il vincitore. Muovendosi nel verso
opposto al gradiente (e assorbendo il 2 in $\eta$) si ottiene la regola.
\end{proof}

È l'algoritmo \termine{K-means} nella versione on-line: una discesa del gradiente stocastica, quindi soggetta a
\textbf{minimi locali} e dipendente dall'\textbf{inizializzazione}; il \textbf{numero di cluster} va fornito in
anticipo.

\subsection{Versione batch}

\begin{algoritmo}[K-means batch]
\begin{enumerate}
  \item Scegliere $K$ centri: $K$ esempi estratti a caso, o $K$ punti casuali nella regione che contiene i dati.
  \item Assegnare ogni esempio al centro più vicino (il \textbf{vincitore}):
        \[
        i^*(\x) = \argmin_i\norm{\x-\w_i}^2,\qquad \norm{\x-\w_i}^2 = \sum_{j=1}^{n}(x_j-w_{ij})^2 .
        \]
  \item Ricalcolare ogni centro come \textbf{media} (centroide) dei membri del suo cluster:
        \[
        \w_i = \frac{1}{\abs{\mathrm{cluster}_i}}\sum_{\x_j\in\mathrm{cluster}_i}\x_j .
        \]
  \item Se il criterio di convergenza non è soddisfatto (nessuna o minima riassegnazione, minima diminuzione dell'errore
        quadratico), tornare al passo 2.
\end{enumerate}
\end{algoritmo}

\begin{proof}[Dalla versione on-line alla batch]
Sommando su tutti gli esempi $\x$ di un cluster $C$ il passo on-line con $\eta=1$ e mediando:
\[
\w_{\mathit{new}} = \w_{\mathit{old}} + \frac{1}{\abs{C}}\sum_{\x\in C}(\x-\w_{\mathit{old}}) = \w_{\mathit{old}} + \frac{1}{\abs C}\sum_{\x\in C}\x - \w_{\mathit{old}} = \frac{1}{\abs C}\sum_{\x\in C}\x ,
\]
cioè la media del cluster, il passo 3 dell'algoritmo batch.
\end{proof}

\begin{figure}[ht]
\centering
\begin{tikzpicture}
\begin{groupplot}[group style={group size=4 by 1, horizontal sep=0.25cm}, scatola, width=4.3cm, height=4.3cm,
  xmin=0, xmax=4, ymin=0, ymax=4, xtick=\empty, ytick=\empty]
\nextgroupplot[title={\footnotesize 1. inizializzazione}]
\addplot[only marks, mark=*, mark size=1pt, grigio] table {dati/km_punti.dat};
\addplot[only marks, mark=x, mark size=4pt, very thick, inchiostro] table {dati/km_centri0.dat};
\nextgroupplot[title={\footnotesize 2. assegnazione}]
\addplot[only marks, mark=*, mark size=1pt, rosso] table[x=x, y=y, restrict expr to domain={\thisrow{c}}{0:0}] {dati/km_ass0.dat};
\addplot[only marks, mark=*, mark size=1pt, blu] table[x=x, y=y, restrict expr to domain={\thisrow{c}}{1:1}] {dati/km_ass0.dat};
\addplot[only marks, mark=*, mark size=1pt, verde] table[x=x, y=y, restrict expr to domain={\thisrow{c}}{2:2}] {dati/km_ass0.dat};
\addplot[only marks, mark=x, mark size=4pt, very thick, inchiostro] table {dati/km_centri0.dat};
\nextgroupplot[title={\footnotesize 3. nuovi centri}]
\addplot[only marks, mark=*, mark size=1pt, rosso] table[x=x, y=y, restrict expr to domain={\thisrow{c}}{0:0}] {dati/km_ass0.dat};
\addplot[only marks, mark=*, mark size=1pt, blu] table[x=x, y=y, restrict expr to domain={\thisrow{c}}{1:1}] {dati/km_ass0.dat};
\addplot[only marks, mark=*, mark size=1pt, verde] table[x=x, y=y, restrict expr to domain={\thisrow{c}}{2:2}] {dati/km_ass0.dat};
\addplot[only marks, mark=x, mark size=4pt, very thick, inchiostro] table {dati/km_centri1.dat};
\nextgroupplot[title={\footnotesize 4. nuova assegnazione}]
\addplot[only marks, mark=*, mark size=1pt, rosso] table[x=x, y=y, restrict expr to domain={\thisrow{c}}{0:0}] {dati/km_ass1.dat};
\addplot[only marks, mark=*, mark size=1pt, blu] table[x=x, y=y, restrict expr to domain={\thisrow{c}}{1:1}] {dati/km_ass1.dat};
\addplot[only marks, mark=*, mark size=1pt, verde] table[x=x, y=y, restrict expr to domain={\thisrow{c}}{2:2}] {dati/km_ass1.dat};
\addplot[only marks, mark=x, mark size=4pt, very thick, inchiostro] table {dati/km_centri1.dat};
\end{groupplot}
\end{tikzpicture}
\caption{K-means batch con $K=3$ (calcolato): inizializzazione dei centri ($\times$), assegnazione degli esempi al centro
più vicino, ricalcolo dei centri come medie, nuova assegnazione.}
\label{fig:kmeans}
\end{figure}

\subsection{Pregi e difetti}

\begin{itemize}
  \item È l'algoritmo più semplice e usato con un criterio d'errore quadratico: facile da implementare ed efficiente.
  \item Il numero $K$ di cluster va fornito.
  \item I minimi locali di $E$ rendono il risultato dipendente dall'inizializzazione, con prototipi solo sub-ottimi.
  \item Funziona molto bene per cluster \textbf{compatti e ipersferici}; con cluster di altre forme (ad esempio due
        spirali intrecciate) fallisce.
  \item \textbf{Nessuna visualizzazione}: non proietta i dati in uno spazio a dimensione minore, e l'indicizzazione dei
        prototipi è arbitraria, quindi la mappatura non è ordinata.
\end{itemize}

Per ridurre il problema dei minimi locali si usa una regola \textbf{soft-max}: si modificano \textbf{tutti} i centri,
con un passo che decresce con la loro distanza da $\x$ (ad esempio il \ingl{maximum-entropy clustering}). Le SOM sono
un'istanza di questa strategia, in cui però la vicinanza tra prototipi è definita su una \textbf{mappa ordinata}.

\section{Self-Organizing Map}

Le \termine{Self-Organizing Map} (SOM, mappe di Kohonen) sono reti di neuroni disposti su una \textbf{griglia regolare
a bassa dimensione} (di solito 2D): ogni unità è identificata dalle sue \textbf{coordinate} $\vv r_i$ sulla mappa,
riceve \textbf{lo stesso input} $\x$ e ha un vettore di pesi $\w_i$ della stessa dimensione $n$ dell'input
(\figrif{fig:som-griglia}).

\begin{figure}[ht]
\centering
\begin{tikzpicture}[x=1cm, y=1cm]
\pgfmathsetseed{4}
\foreach \k in {1,...,40} \fill[grigio] ({0.4+2.2*rnd},{0.4+2.2*rnd}) circle (1.1pt);
\draw[grigio] (0,0) rectangle (3,3);
\node[font=\footnotesize] at (1.5,-0.35) {spazio di input ($\R^n$)};
\fill[rosso] (1.9,2.1) circle (2pt) node[above, font=\scriptsize] {$\x$};
\begin{scope}[xshift=5.5cm, yshift=0.2cm]
\foreach \i in {0,...,5} \foreach \j in {0,...,4} \node[circle, draw, inner sep=0pt, minimum size=3.2mm] (u\i\j) at (\i*0.5,\j*0.55) {};
\foreach \i in {0,...,5} \foreach \j in {0,...,3} { \pgfmathtruncatemacro{\jj}{\j+1} \draw[grigio] (u\i\j) -- (u\i\jj); }
\foreach \i in {0,...,4} \foreach \j in {0,...,4} { \pgfmathtruncatemacro{\ii}{\i+1} \draw[grigio] (u\i\j) -- (u\ii\j); }
\fill[rosso] (u32) circle (1.6mm);
\foreach \a/\b in {2/2,4/2,3/1,3/3} \fill[rosso!40] (u\a\b) circle (1.6mm);
\node[font=\footnotesize, align=center] at (1.25,-0.6) {mappa 2D di unità};
\end{scope}
\draw[->, rosso, thick] (2.0,2.2) to[bend left=15] node[above, font=\scriptsize] {vincitore} (6.9,1.35);
\end{tikzpicture}
\caption{Una SOM: ogni unità della griglia ha un vettore di pesi nello spazio di input. Il vincitore è l'unità con i pesi
più vicini all'input; i suoi vicini sulla mappa (rosa) vengono aggiornati con lui.}
\label{fig:som-griglia}
\end{figure}

La SOM apprende una mappa dallo spazio di input alla griglia che \textbf{preserva la topologia}: unità vicine sulla
mappa rispondono a input simili, e punti vicini nello spazio di input vengono mappati sulla stessa unità o su unità
vicine. Si usa per il \textbf{clustering} (i cluster si individuano sulla mappa), per la \textbf{VQ} (l'input diventa il
vettore dei pesi del vincitore), per la \textbf{compressione} (l'input diventa l'indice del vincitore) e per la
\textbf{proiezione ed esplorazione} dei dati in due dimensioni.

L'apprendimento è \termine{competitivo}: i neuroni competono per ogni input, il vincitore apprende di più, e
gradualmente le unità si specializzano su insiemi diversi di input. Un analogo biologico è la corteccia
somatosensoriale, in cui una mappa topologicamente ordinata del corpo associa neuroni adiacenti a parti del corpo
vicine.

\begin{algoritmo}[SOM]
\begin{enumerate}
  \item Inizializzare casualmente i pesi della mappa.
  \item Estrarre un input $\x$.
  \item \textbf{Fase competitiva}: trovare il vincitore, l'unità con $\w$ più simile a $\x$ (\ingl{winner-take-all}):
        \[
        i^*(\x) = \argmin_i\norm{\x-\w_i}.
        \]
  \item \textbf{Fase cooperativa}: aggiornare i pesi del vincitore \textbf{e delle unità a lui vicine sulla mappa}:
        \begin{equation}\label{eq:som}
        \w_i(t+1) = \w_i(t) + \eta(t)\,h_{i,i^*(\x)}(t)\,\big[\x-\w_i(t)\big].
        \end{equation}
  \item Ripetere fino a convergenza.
\end{enumerate}
\end{algoritmo}

La fase cooperativa è un aggiornamento hebbiano (i pesi si avvicinano all'input) modulato dalla \termine{funzione di
vicinato} $h$, che decresce con la distanza \textbf{sulla mappa} tra l'unità $i$ e il vincitore, ad esempio una
gaussiana
\[
h_{i,i^*}(t) = \exp\left(-\frac{\norm{\vv r_i-\vv r_{i^*}}^2}{2\sigma^2(t)}\right),
\]
con $\sigma(t)$ larghezza del vicinato (in pratica si usano anche vicinati di forma geometrica, rettangolare o
esagonale, con un numero finito di unità). Il \textbf{learning rate} $\eta$ e il \textbf{raggio} $\sigma$ diminuiscono
con le iterazioni: il vicinato è \textbf{ampio all'inizio} (ordinamento globale della mappa) e si restringe lentamente
(convergenza e raffinamento locale).

\begin{idea}[L'emergere dell'ordine topologico]
I pesi non sono modificati indipendentemente ma \textbf{per sottoinsiemi topologicamente legati}: il vincitore e i suoi
vicini sulla griglia ricevono aggiornamenti simili e quindi, dopo l'apprendimento, rispondono a input simili. Così
l'informazione topologica viene trasmessa alla mappa (\figrif{fig:som-training}). Una dimostrazione formale
dell'ordinamento è però difficile.
\end{idea}

\begin{figure}[ht]
\centering
\begin{tikzpicture}
\begin{groupplot}[group style={group size=6 by 1, horizontal sep=0.18cm}, scatola, width=3.45cm, height=3.45cm,
  xmin=-0.05, xmax=1.05, ymin=-0.05, ymax=1.05, xtick=\empty, ytick=\empty]
\nextgroupplot[title={\scriptsize $t=0$}]
\addplot[thin, blu] table {dati/som_0.dat};
\nextgroupplot[title={\scriptsize $t=20$}]
\addplot[thin, blu] table {dati/som_20.dat};
\nextgroupplot[title={\scriptsize $t=100$}]
\addplot[thin, blu] table {dati/som_100.dat};
\nextgroupplot[title={\scriptsize $t=1000$}]
\addplot[thin, blu] table {dati/som_1000.dat};
\nextgroupplot[title={\scriptsize $t=5000$}]
\addplot[thin, blu] table {dati/som_5000.dat};
\nextgroupplot[title={\scriptsize $t=20000$}]
\addplot[thin, blu] table {dati/som_20000.dat};
\end{groupplot}
\end{tikzpicture}
\caption{Addestramento di una SOM $10\times10$ su dati uniformi nel quadrato (calcolato): i vettori dei pesi, collegati
secondo la griglia, partono concentrati al centro e si distendono progressivamente sulla distribuzione.}
\label{fig:som-training}
\end{figure}

\begin{table}[ht]
\centering\small
\caption{K-means e SOM a confronto.}
\begin{tabularx}{0.95\linewidth}{XX}
\toprule
K-means & SOM \\
\midrule
aggiorna solo il vincitore & aggiorna il vincitore e i suoi vicini sulla griglia, con un vicinato che si restringe \\
prototipi senza alcun ordine & prototipi su una mappa ordinata che preserva la topologia \\
nessuna visualizzazione & visualizzazione 2D dei dati \\
sensibile ai minimi locali & il vicinato ampio iniziale aiuta a evitarli \\
\bottomrule
\end{tabularx}
\end{table}
Con un vicinato di raggio zero la SOM si riduce al K-means on-line.

\subsection{Visualizzazione}

Sulla mappa si possono visualizzare la \textbf{densità} dei vettori di riferimento (colore proporzionale al numero di
esempi assegnati a ogni unità) e le \textbf{distanze} tra i pesi di unità vicine, con la \termine{U-matrix}: colori scuri
indicano grandi distanze tra unità adiacenti (confini tra cluster), colori chiari piccole distanze (unità dello stesso
cluster). Questo rende la SOM facilmente interpretabile. Ad esempio, addestrando una mappa $13\times9$ su 39 indicatori
di benessere di 77 paesi, le zone chiare della U-matrix raggruppano paesi simili e le zone scure separano i gruppi,
senza alcuna ipotesi a priori sulla forma dei cluster.

\section{Valutare un clustering}

La valutazione è spesso \textbf{soggettiva}: esistono pochi riferimenti oggettivi, tranne in sotto-domini ben definiti.
Esistono misure \textbf{oggettive} come l'errore di quantizzazione, che però dipende dal numero di centroidi. Non si usa
semplicemente l'etichetta di un task di classificazione sottostante (la \emph{purity}): se si hanno le etichette, tanto
vale usare un classificatore. Va distinta la qualità del clustering in sé dalla qualità dei risultati per il dominio.
Altri metodi non supervisionati sono il clustering gerarchico (dendrogrammi), l'analisi delle componenti principali, le
regole di associazione, ICA, Multidimensional Scaling, t-SNE.

\begin{puntichiave}
\begin{itemize}
  \item VQ: codifica i dati con un codebook; loss = errore di quantizzazione $\sum\norm{\x_i-\w_{i^*}}^2$.
  \item K-means on-line: $\Delta\w_{i^*} = \eta(\x-\w_{i^*})$ solo per il vincitore; batch: assegnazione e media.
  \item SOM: griglia di unità; fase competitiva (vincitore) e cooperativa (vincitore e vicini sulla mappa, con vicinato
        gaussiano che si restringe); preserva la topologia e permette la visualizzazione (U-matrix).
\end{itemize}
\end{puntichiave}

\begin{domande}
\begin{enumerate}
  \item Definire clustering e quantizzazione vettoriale. Cos'è l'errore di quantizzazione?
  \item Derivare la regola on-line del K-means e passare alla versione batch. Pregi e difetti del K-means.
  \item Descrivere la SOM: architettura, fase competitiva, fase cooperativa, funzione di vicinato.
  \item Perché la SOM preserva la topologia? Differenze tra K-means e SOM.
  \item Come si visualizza una SOM?
\end{enumerate}
\end{domande}
